\chapter{Mengalikan Bilangan Negatif}\label{ch:g8:negprod}

Kelas 7 menjumlahkan dan mengurangkan \omterm{def:g7:negatives:def}{bilangan bertanda}
(\cref{ch:g7:negatives}); yang tersisa adalah mengalikan dan
membaginya. Satu pertanyaan menguasai bab ini: mengapa ``minus kali
minus'' harus menjadi plus? Kita memberi alasannya, bukan sekadar
aturannya.

\section{Aturan tanda}

\begin{theorem}[Tanda sebuah hasil kali]\label{thm:g8:negprod:rules}
\omterm{def:g2:mult:def}{Hasil kali} dua \omterm{def:g7:negatives:def}{bilangan bertanda} punya jarak ke \omterm{def:g1:counting:zero}{nol} yang sama dengan
\omterm{def:g2:mult:def}{hasil kali} kedua jaraknya, dan tandanya diberikan oleh:
\begin{center}
\renewcommand{\arraystretch}{1.3}
\begin{tabular}{c|cc}
$\times$ & $+$ & $-$ \\
\hline
$+$ & $+$ & $-$ \\
$-$ & $-$ & $+$ \\
\end{tabular}
\end{center}
\emph{Tanda yang sama: \omterm{def:g2:mult:def}{hasil kalinya} positif. Tanda yang berlawanan:
\omterm{def:g2:mult:def}{hasil kalinya} negatif.} Tabel yang sama berlaku untuk \omterm{def:g3:division:remainder}{hasil bagi}.
\end{theorem}

\begin{proof}[Mengapa $(-1) \times (-1) = +1$]
Pertama, $3 \times (-4)$ berarti $(-4) + (-4) + (-4) = -12$: positif
kali negatif adalah negatif. Sekarang perhatikan polanya:
\[
3 \times (-4) = -12, \quad
2 \times (-4) = -8, \quad
1 \times (-4) = -4, \quad
0 \times (-4) = 0 .
\]
Pada tiap langkahnya, faktor pertamanya turun $1$ dan hasilnya
\emph{naik} $4$. Dengan melanjutkan polanya satu langkah lagi:
\[
(-1) \times (-4) = 0 + 4 = +4 .
\]
Pilihan lain apa pun akan merusak keteraturan aritmetika (yaitu sifat
distributif). Jadi minus kali minus haruslah plus.
\end{proof}

\begin{omfigure}
\begin{tikzpicture}
\begin{axis}[omaxis,
  width=7.8cm, height=5.4cm,
  xmin=-2.6, xmax=3.6, ymin=-14, ymax=10,
  xlabel={$k$}, ylabel={$k \times (-4)$},
  xtick={-2,-1,1,2,3}, ytick={-12,-8,-4,4,8}]
  \addplot[omDef, thick, domain=-2.3:3.3] {-4*x};
  \addplot[omDef, only marks, mark=*, mark size=1.6pt] coordinates
    {(3,-12) (2,-8) (1,-4) (0,0) (-1,4) (-2,8)};
\end{axis}
\end{tikzpicture}

{\small Alasan polanya dalam satu gambar: titik
$(k,\ k \times (-4))$ berjajar, dan melanjutkan garisnya melewati
$k = 0$ memaksa $(-1)\times(-4) = +4$ dan $(-2)\times(-4) = +8$.}
\end{omfigure}

\begin{example}\label{ex:g8:negprod:products}
\begin{align*}
(-7) \times (+6) &= -42 && \text{(tanda berlawanan)} \\
(-5) \times (-8) &= +40 && \text{(tanda sama)} \\
(+2.5) \times (-4) &= -10 \\
(-36) \div (-9) &= +4 \\
(+15) \div (-2) &= -7.5 .
\end{align*}
\end{example}

\section{Hasil kali beberapa faktor}

\begin{proposition}[Tanda hasil kali yang panjang]\label{prop:g8:negprod:many}
Tanda \omterm{def:g2:mult:def}{hasil kali} beberapa faktor yang bukan \omterm{def:g1:counting:zero}{nol} hanya bergantung pada
banyaknya faktor yang negatif:
\begin{itemize}
  \item banyaknya faktor negatif yang \emph{\omterm{def:g3:numbers:evenodd}{genap}} memberi \omterm{def:g2:mult:def}{hasil kali}
        yang positif;
  \item banyaknya yang \emph{\omterm{def:g3:numbers:evenodd}{ganjil}} memberi \omterm{def:g2:mult:def}{hasil kali} yang negatif.
\end{itemize}
\end{proposition}

\begin{proof}
Faktor negatifnya berpasangan, dan tiap pasangan menyumbang tanda
positif (\cref{thm:g8:negprod:rules}); banyaknya yang \omterm{def:g3:numbers:evenodd}{ganjil}
menyisakan satu faktor negatif tanpa pasangan, yang membuat seluruh
\omterm{def:g2:mult:def}{hasil kalinya} negatif.
\end{proof}

\begin{example}\label{ex:g8:negprod:many}
$(-2) \times (+3) \times (-5) \times (-1)$: tiga faktor negatif
(\omterm{def:g3:numbers:evenodd}{ganjil}), jadi \omterm{def:g2:mult:def}{hasil kalinya} negatif; jaraknya: $2 \times 3 \times 5
\times 1 = 30$. Hasilnya: $-30$. Tentukan tandanya \emph{lebih
dahulu}, lalu kalikan jaraknya --- dua tugas yang mudah, bukan satu
tugas yang rawan galat.
\end{example}

\begin{example}[Pangkat bilangan negatif]\label{ex:g8:negprod:powers}
$(-2)^2 = (-2) \times (-2) = +4$, dan
$(-2)^3 = (-2)^2 \times (-2) = -8$: pangkat \omterm{def:g3:numbers:evenodd}{genap} memberi nilai yang
positif, pangkat \omterm{def:g3:numbers:evenodd}{ganjil} mempertahankan tandanya. Hati-hati dengan
notasinya: $(-3)^2 = 9$ tetapi $-3^2 = -(3 \times 3) = -9$.
\end{example}

\section{Menghitung dengan keempat operasinya}

\begin{method}[Menghitung dengan aman memakai bilangan bertanda]\label{met:g8:negprod:compute}
\begin{enumerate}
  \item Patuhi prioritas pada \cref{ch:g7:priorities}: tanda kurung,
        lalu $\times$ dan $\div$, lalu $+$ dan $-$;
  \item pada tiap perkalian atau pembagian, carilah \emph{tandanya
        lebih dahulu}, lalu jaraknya;
  \item tulislah ulang satu langkah pada tiap barisnya.
\end{enumerate}
\end{method}

\begin{example}\label{ex:g8:negprod:compute}
\begin{align*}
5 - 3 \times (-4)
&= 5 - (-12) && \text{(hasil kalinya dahulu: tanda berlawanan)}\\
&= 5 + 12 = 17 .
\end{align*}
\begin{align*}
\frac{(-6) \times (-10)}{-4}
&= \frac{+60}{-4} && \text{(pembilangnya: tanda sama)}\\
&= -15 && \text{(hasil baginya: tanda berlawanan).}
\end{align*}
\end{example}

\begin{example}[Menyulih nilai yang negatif]\label{ex:g8:negprod:substitute}
Hitunglah $E = 3x^2 - 2x$ untuk $x = -5$, dengan tanda kurung mengelilingi nilainya:
\[
E = 3 \times (-5)^2 - 2 \times (-5)
= 3 \times 25 - (-10)
= 75 + 10 = 85 .
\]
Tanda kurung pada $(-5)^2$ sangat penting: tanpanya kuadratnya hanya
akan berlaku untuk $5$ saja.
\end{example}

\section{Latihan}

\begin{exercise}[$\star$]\label{exo:g8:negprod:1}
Hitunglah:
\[
(-8) \times (+7), \qquad
(-9) \times (-6), \qquad
(+12) \times (-0.5), \qquad
(-1) \times (-1) \times (-1).
\]
\end{exercise}

\begin{exercise}[$\star$]\label{exo:g8:negprod:2}
Hitunglah:
\[
(-63) \div (+9), \qquad
(-8) \div (-16), \qquad
\frac{45}{-5}, \qquad
\frac{-7.2}{-8} .
\]
\end{exercise}

\begin{exercise}[$\star$]\label{exo:g8:negprod:3}
Sebutkan hanya \emph{tanda} tiap \omterm{def:g2:mult:def}{hasil kalinya}, tanpa menghitungnya:
\[
(-3) \times (-7) \times (+2) \times (-5); \qquad
(-1)^{10}; \qquad
(-2)^{7}; \qquad
(-4) \times 0 \times (-6).
\]
\end{exercise}

\begin{exercise}[$\star$]\label{exo:g8:negprod:4}
Hitunglah:
\[
(-2)^4, \qquad
-2^4, \qquad
(-3)^3, \qquad
(-1)^{2026} .
\]
\end{exercise}

\begin{exercise}[$\star$]\label{exo:g8:negprod:5}
Hitunglah selangkah demi selangkah, dengan mematuhi prioritasnya:
\[
7 + 2 \times (-6), \qquad
(-4) \times (5 - 9), \qquad
18 \div (-3) - (-2) .
\]
\end{exercise}

\begin{exercise}[$\star$]\label{exo:g8:negprod:6}
Hitunglah:
\[
\frac{(-5) \times (+8)}{-4}, \qquad
\frac{(-3) \times (-6)}{(-2) \times (+9)} .
\]
\end{exercise}

\begin{exercise}[$\star$]\label{exo:g8:negprod:7}
Hitunglah untuk $x = -3$:
\[
5x, \qquad x^2, \qquad -x^2, \qquad 2x^2 + 4x, \qquad (2x)^2 .
\]
\end{exercise}

\begin{exercise}[$\star\star$]\label{exo:g8:negprod:8}
Lengkapilah tiap kesamaannya:
\[
(-6) \times \;?\; = 42, \qquad
\;?\; \div (-4) = -2.5, \qquad
(-5) \times \;?\; = -1 .
\]
\end{exercise}

\begin{exercise}[$\star\star$]\label{exo:g8:negprod:9}
Benar atau salah? Berilah alasan atau contoh penyangkalnya.
\begin{enumerate}
  \item \omterm{def:g2:mult:def}{Hasil kali} dua \omterm{def:g7:negatives:def}{bilangan bertanda} selalu sedikitnya sama besar
        dengan \omterm{def:g1:addition:def}{jumlahnya}.
  \item Kuadrat sebuah \omterm{def:g7:negatives:def}{bilangan bertanda} tidak pernah negatif.
  \item Kalau \omterm{def:g2:mult:def}{hasil kali} tiga faktor bernilai positif, ketiga
        faktornya positif.
\end{enumerate}
\end{exercise}

\begin{exercise}[$\star\star$]\label{exo:g8:negprod:10}
Tiap jawaban yang salah pada sebuah kuis bernilai $-3$ angka, tiap
jawaban yang benar $+5$. Zahra menjawab $20$ pertanyaan dan
memperoleh $36$ angka. Berapa jawabannya yang benar? (Cobalah
beberapa nilai, atau susunlah perhitungan
$5r - 3(20 - r) = 36$.)
\end{exercise}

\begin{exercise}[$\star\star\star$]\label{exo:g8:negprod:11}
Dengan memakai sifat distributif (seperti pada bukti
\cref{thm:g8:negprod:rules}), jabarkan dan berilah alasan untuk tiap
langkahnya pada
\[
0 = (-1) \times 0 = (-1) \times \bigl(1 + (-1)\bigr)
= (-1) \times 1 + (-1) \times (-1),
\]
lalu simpulkan bahwa $(-1) \times (-1)$ harus sama dengan $+1$.
\end{exercise}

\section{Soal: Mengapa minus kali minus adalah plus}

\begin{problem}[{Soal akhir pekan --- aturan tanda yang disimpulkan
dari sifat distributif, dan jumlah berganti tanda $1 - 2 + 3 - 4 + \dots$}]
\label{pb:g8:negprod:1}
Bukti \cref{thm:g8:negprod:rules} melanjutkan sebuah pola lalu
menyatakan bahwa ``pilihan lain apa pun akan merusak sifat
distributif''. Soal ini membuat pernyataan itu tepat: berangkat hanya
dari sifat distributif
(\cref{thm:g7:priorities:distributivity}), kamu akan
\emph{membuktikan} aturan tandanya --- tanpa gambar, tanpa pola,
dengan cara yang dipakai aljabar --- lalu mempekerjakannya pada
pangkat $-1$ dan pada \omterm{def:g1:addition:def}{jumlah} terkenal yang tandanya berganti-ganti.
Sepanjang soal ini, $a$ dan $b$ adalah \omterm{def:g7:negatives:def}{bilangan bertanda}; kita
menerima begitu saja aturan penjumlahan pada \cref{ch:g7:negatives},
bahwa $1 \times a = a$, dan bahwa \omterm{def:g2:mult:def}{hasil kali} dapat dihitung dalam
urutan mana pun (jadi sifat distributif berlaku pada kedua sisi
sebuah \omterm{def:g2:mult:def}{hasil kali}).

\medskip
\noindent\textbf{Bagian I --- Aturan tandanya adalah teorema.}
\begin{enumerate}
  \item Buktikan bahwa $0 \times a = 0$ untuk setiap \omterm{def:g7:negatives:def}{bilangan
        bertanda} $a$. (Tulislah $0 = 0 + 0$, jabarkan
        $(0 + 0) \times a$, lalu tanyakan bilangan mana yang, kalau
        ditambahkan ke $0 \times a$, memberi $0 \times a$ kembali.)
  \item Jabarkan $\bigl(1 + (-1)\bigr) \times a$ untuk menunjukkan
        bahwa $(-1) \times a + a = 0$, lalu simpulkan bahwa
        \[
        (-1) \times a = -a :
        \]
        mengalikan dengan $-1$ memberi \omterm{def:g7:negatives:def}{lawannya}. Bandingkan dengan
        \cref{exo:g8:negprod:11}, yaitu keadaan $a = -1$.
  \item Simpulkan bahwa $(-a) \times b = -(a \times b)$: satu tanda
        minus keluar dari \omterm{def:g2:mult:def}{hasil kali} tanpa berubah.
  \item Simpulkan bahwa $(-a) \times (-b) = a \times b$: dua tanda
        minus saling meniadakan.
  \item Setiap \omterm{def:g7:negatives:def}{bilangan bertanda} adalah jaraknya ke \omterm{def:g1:counting:zero}{nol} dengan sebuah
        tanda di depannya: $a = +d$ atau $a = -d$. Jelaskan bagaimana
        pertanyaan 3 dan~4 membuktikan keempat keadaan tabel tanda
        pada \cref{thm:g8:negprod:rules}, termasuk jaraknya.
\end{enumerate}

\medskip
\noindent\textbf{Bagian II --- Pangkat $-1$.}
\begin{enumerate}[resume]
  \item Hitunglah $(-1)^2$, $(-1)^3$, $(-1)^4$, $(-1)^5$, lalu
        sebutkan, dengan alasan dari \cref{prop:g8:negprod:many},
        nilai $(-1)^n$ untuk setiap bilangan bulat $n \geq 1$.
  \item Sebutkan (tanpa menghitung satu jarak pun) tanda dari
        \[
        (-1) \times (-2) \times (-3) \times \dots \times (-10),
        \]
        lalu tulislah jarak \omterm{def:g2:mult:def}{hasil kali} ini sebagai \omterm{def:g2:mult:def}{hasil kali}
        bilangan bulat, tanpa mengerjakannya.
  \item Tunjukkan bahwa $(-a)^2 = a^2$ dan bahwa $(-a)^3 = -a^3$.
        Untuk pangkat yang mana tanda minusnya bertahan?
  \item Pakailah aturan tandanya untuk menunjukkan bahwa kuadrat
        sebuah \omterm{def:g7:negatives:def}{bilangan bertanda} tidak pernah negatif, lalu simpulkan
        bahwa tidak ada \omterm{def:g7:negatives:def}{bilangan bertanda} $x$ yang memenuhi
        $x^2 = -4$.
  \item Zahra menyatakan: ``$(-1)^{2026} + (-1)^{2027} = 0$, dan
        lebih umum lagi dua pangkat $-1$ yang berurutan selalu saling
        meniadakan.'' Apakah dia benar? Berilah alasan.
\end{enumerate}

\medskip
\noindent\textbf{Bagian III --- \omterm{def:g1:addition:def}{Jumlah} yang berganti tanda.}
Dengan aturan tandanya yang sudah aman, kita dapat menghitung \omterm{def:g1:addition:def}{jumlah}
yang mencampur kedua tandanya. Untuk bilangan bulat $n \geq 1$,
misalkan
\[
A_n = 1 - 2 + 3 - 4 + \dots
\]
adalah \omterm{def:g1:addition:def}{jumlah} bilangan bulat dari $1$ sampai $n$ dengan tanda yang
berganti-ganti: suku terakhirnya $+n$ ketika $n$ \omterm{def:g3:numbers:evenodd}{ganjil}, dan $-n$
ketika $n$ \omterm{def:g3:numbers:evenodd}{genap}. Jadi $A_1 = 1$ dan $A_2 = 1 - 2 = -1$.
\begin{enumerate}[resume]
  \item Hitunglah $A_3$, $A_4$, $A_5$, $A_6$ dan $A_7$. Apa yang kamu
        duga?
  \item Andaikan $n$ \omterm{def:g3:numbers:evenodd}{genap}. Kelompokkan sukunya dua-dua,
        $(1 - 2) + (3 - 4) + \dots$, bilanglah pasangannya, lalu
        buktikan bahwa $A_n = -\frac{n}{2}$.
  \item Andaikan $n$ \omterm{def:g3:numbers:evenodd}{ganjil}. Jelaskan mengapa $A_n = A_{n-1} + n$,
        lalu simpulkan dari pertanyaan sebelumnya bahwa
        $A_n = \frac{n + 1}{2}$.
  \item Hitunglah $A_{100}$, $A_{101}$ dan $A_{2026}$.
  \item Alma menghitung $A_{101} - A_{100} = 51 - (-50) = 101$ lalu
        terkejut menemukan tepat $101$. Seharusnya dia terkejut?
        Jelaskan dalam satu kalimat, lalu sebutkan nilai
        $A_{2027} - A_{2026}$ tanpa perhitungan lebih lanjut.

\end{enumerate}
\end{problem}
